\begin{textsample}{POS Dim 4 – human – Score 5.00 – mg/mg\_ef\_ciencias\_da\_natureza\_intro\_p3.txt}  \label{ex:f4_pos_014}
Em Ciências da Natureza, os procedimentos que correspondem aos modos de buscar e organizar conhecimentos são bastante variados : a observação, a experimentação, a comparação, a elaboração de hipóteses e suposições, o debate \textbf{oral} sobre hipóteses, o estabelecimento de relações entre fatos ou fenômenos e ideias por meio da \textbf{leitura} e \textbf{escrita} de \textbf{textos} informativos, a elaboração de roteiros de pesquisa bibliográfica e questões para enquete, a busca de informações em fontes variadas, a organização de informações por meio de desenhos, tabelas, gráficos, esquemas e \textbf{textos}, o confronto entre suposições entre elas e os dados obtidos por investigação, a elaboração de perguntas e problemas, a proposição para a solução de problemas.
Pensando assim, o ensino de Ciências deve promover situações nas quais os estudantes possam :

% matched lemmas: escrita, leitura, oral, texto
\end{textsample}
